\subsection{Software Tools}

Such a heterogeneous combination of development platforms, libraries, and simulators is applied for developing the system and validating its efficiency. In particular, the chosen set of software components ensures compatibility of the system with ESP32 microcontroller and supports implementation of real-time signal processing and reinforcement learning algorithm.

\subsubsection{Arduino IDE}
The Arduino IDE is applied for programming and flashing the firmware to the ESP32 DevKit microcontroller. The integrated development platform features a user-friendly interface, built-in ESP32 libraries, and easy-to-use peripherals support for working with pulse-width modulation (PWM) motors, analog-to-digital conversion (ADC) sensors, and Wi-Fi modules. With Arduino IDE, developers can prototype their algorithms, test the functionality of RSSI measurement and servo motor movement, and perform serial communication debugging.

\subsubsection{Programming Languages}

\textbf{Python:} The Python programming language is used for conducting simulations, analysis of collected datasets, and reinforcement learning experiment. Namely, such libraries as NumPy, Matplotlib, and Pandas allow performing signal processing operations, analyzing visualized RSSI results, and testing reinforcement learning policies before their deployment onto microcontrollers. 

Additionally, Python allows conducting post-processing of collected data and analyzing the performance of the proposed solution in different environmental conditions. Besides that, Python allows designing a simulation script modeling antenna orientation, RSSI feedback, and reinforcement learning policies. Off-line validation of lightweight tabular reinforcement learning policies will be conducted using Python scripts and their further translation into on-board logic for ESP32.

\textbf{C/C++:} The firmware running on ESP32 microcontroller will be written in C/C++ using Arduino IDE/PlatformIO toolchains. Such a choice is motivated by the need to have deterministic execution, low-level hardware support, and efficient memory management when implementing complex algorithms like RSSI measurement and servo motor control.

\subsubsection{PlatformIO}

PlatformIO is an advanced firmware development environment for ESP32 microcontrollers based on Visual Studio Code. It offers better support for dependencies management, multi-environment building, and debugging compared to the Arduino IDE.

\subsubsection{Version Control}

Git and GitHub services are applied for maintaining version histories of all source code, configurations, and documentation developed within the course of the research.

\subsubsection{Trello}

Trello is used for managing and tracking project-related tasks assigned to different members of the research team. The described application allows organizing tasks in boards, lists, and cards in order to facilitate project management.

\subsubsection{Draw.io}

The Draw.io diagrams editor is used for creating system diagrams and flowcharts. It provides users with a convenient way of making structured and understandable visuals.

\subsubsection{KiCad}

KiCad EDA was used to create the schematics and PCB layout for the design. It was helpful in arranging the components and customizing the power and signal lines, as well as ensuring slip rings are connected. The three-dimensional feature provided by KiCad was useful in confirming mechanical clearance and ensuring that the moving and stationary parts are aligned correctly.